\begin{figure}
  \centering
  \begin{tikzpicture}
    \begin{scope}
      \clip[rounded corners=9pt]
      (-0.24\textwidth,-0.24\textwidth) rectangle
      (0.24\textwidth,0.24\textwidth);
      \node[inner sep=0pt,outer sep=0pt] (art) {
      \includegraphics[width=0.48\textwidth]{media/guitar-sandstorm/illustration.png}};
      % Soft white rings blend the image into the paper without changing the art.
      \begin{scope}
        \clip (art.south west) rectangle (art.north east);
        \foreach \k in {24,23,...,1} {
        \draw[white,opacity=0.12,line width=\k*0.65pt,
        rounded corners=9pt]
        (art.south west) rectangle (art.north east);
        }
      \end{scope}
      % The outer white frame also removes the square corner remnants.
      \draw[white,line width=3pt,rounded corners=9pt]
      (art.south west) rectangle (art.north east);
    \end{scope}
  \end{tikzpicture}
  \caption[A guitar in a sandstorm]{AI-generated illustration inspired by
    Walsh's guitar-in-a-sandstorm example. Artwork from Le Chen's
    \href{https://github.com/SPDEngine/SPDE-Visuals}{SPDE-Visuals collection},
    shared under the
    \href{https://creativecommons.org/publicdomain/zero/1.0/}{CC0 public-domain dedication}.
    This is an artistic illustration, not a
    numerical simulation. Walsh's worked example uses an infinite string to
    avoid boundary conditions. Companion videos and source information are
    included in \path{media/guitar-sandstorm/}.}
  \label{Fig:walsh-guitar}
\end{figure}
